\documentclass[10pt,twocolumn]{article}

\usepackage[margin=0.72in]{geometry}
\usepackage{booktabs}
\usepackage{microtype}
\usepackage{siunitx}
\usepackage{tabularx}
\usepackage{xcolor}
\usepackage{hyperref}

\sisetup{detect-all,round-mode=places,round-precision=3}
\newcommand{\hpeek}{H-PEEK}
\newcommand{\peek}{PEEK}
\newcommand{\todoresult}{\textcolor{gray}{pending}}

\title{H-PEEK: History-Enhanced Predictive KV Cache Management\\
       \large Prototype Design and Preliminary Evaluation}
\author{H-PEEK Project}
\date{October 19, 2026}

\begin{document}
\maketitle

\begin{abstract}
PEEK protects reusable KV-cache blocks using demand visible in the current
waiting queue. This report evaluates a complementary cross-request signal: the
frequency with which a prefix block appeared in a bounded window of previously
admitted requests. We implement this signal in the PEEK Online integration for
vLLM V1, without replacing PEEK's scheduler. On a deterministic periodic-reuse
trace, H-PEEK preserves an absent 2,992-token prefix across cache pressure and
reduces its return-request time-to-first-token (TTFT) from 272.5 to 50.9 ms.
Across three seeds, throughput changes by +2.57\% on the positive trace,
-0.10\% on a no-sharing trace, and -0.80\% on PEEK's concurrent queue-sharing
case. Internal eviction time and scan volume increase substantially, motivating
the overhead and adaptation ablations specified in this report. These are
prototype results, not a reproduction of all PEEK paper results.
\end{abstract}

\section{Scope and Status}

The current artifact validates the H-PEEK eviction mechanism on one NVIDIA RTX
A5000 and Qwen2.5-3B-Instruct. It does not yet claim full reproduction of the
PEEK paper or final 7B performance. All reported paired comparisons use the
same PEEK cLPM scheduler; only the eviction policy changes.

Implementation validation includes 195 passing upstream PEEK Python CPU tests,
24 passing Rust tests, 59 passing online regression tests (six skipped), and
nine H-PEEK-specific history/BlockPool tests. The required zero-waiting-demand
case, exact window expiry, empty-history fallback, key deduplication, memory
bound, and KV-cache-group identity are covered.

\begin{table}[t]
  \centering
  \caption{Locked prototype environment.}
  \label{tab:environment}
  \begin{tabular}{ll}
    \toprule
    Component & Version or setting \\
    \midrule
    PEEK upstream & \texttt{3aca72b63c1d} \\
    vLLM & 0.19.1, V1 engine \\
    PyTorch & 2.10.0+cu128 \\
    Model & Qwen2.5-3B-Instruct \\
    GPU & 1$\times$ RTX A5000 24 GB \\
    KV capacity & 4,096 blocks / 65,536 tokens \\
    History window & 256 admitted requests \\
    History weight $\alpha$ & 1.0 \\
    Seeds & 42, 142, 242 \\
    \bottomrule
  \end{tabular}
\end{table}

\section{Method}

\subsection{Admitted-request history}

Each independent request contributes at most one observation when it first
moves from \texttt{WAITING} to \texttt{RUNNING}. Preemption, resume, and
re-entry into the waiting queue do not add observations. Prefix hashes are
deduplicated within a request and stored in an exact FIFO request window. When
the oldest request expires, its contribution is decremented. Consequently,
retained references are bounded by
$W\times K_{\max}$, where $W$ is the request-window size and $K_{\max}$ is the
per-request key cap.

For an evictable block $b$, the MVP protection score is
\begin{equation}
  S(b) = q_b + \alpha h_b,
  \label{eq:score}
\end{equation}
where $q_b$ is current waiting-queue demand and $h_b$ is the number of admitted
requests in the active window containing $b$. Higher scores are harder to
evict. This request-frequency signal is distinct from PEEK's block access-time
\texttt{recency} and \texttt{decay} modes.

\subsection{Pipeline integration}

The implementation makes two behaviorally critical changes. First, the
eviction path falls back to stock LRU only when neither usable waiting demand
nor usable historical demand exists. Second, a candidate with $q_b=0$ is still
scored using $h_b$ rather than being selected immediately. Hash identity uses
vLLM's \texttt{BlockHashWithGroupId} conversion so cache groups do not alias.
The existing free queue, reference counts, cached-hash mapping, and scan cap
remain authoritative.

\begin{table*}[t]
  \centering
  \caption{Policy definitions used in the prototype.}
  \label{tab:policies}
  \begin{tabularx}{\textwidth}{lllX}
    \toprule
    Policy & Scheduler & Eviction signal & Purpose \\
    \midrule
    Vanilla & vLLM & LRU/APC & Engine baseline \\
    PEEK cluster & PEEK cLPM & Current queue demand and prefix depth & Primary PEEK baseline \\
    PEEK recency & PEEK-compatible & Queue demand and linear access age & Existing age baseline \\
    PEEK decay & PEEK-compatible & Queue demand and exponential access age & Existing age baseline \\
    H-PEEK history-only & vLLM/PEEK & $\alpha h_b$ & History ablation \\
    H-PEEK joint & PEEK cLPM & $q_b+\alpha h_b$ & Proposed P1 policy \\
    \bottomrule
  \end{tabularx}
\end{table*}

\section{Experimental Design}

\begin{table*}[t]
  \centering
  \caption{Completed and planned workloads. Cache state is cleared before each
  measured workload.}
  \label{tab:workloads}
  \begin{tabularx}{\textwidth}{lllX}
    \toprule
    Workload & Arrival mode & Status & Property under test \\
    \midrule
    Periodic reuse & Sequential & Complete, 3 seeds & Prefix absent from waiting queue returns after pressure \\
    Single-target drift & Sequential & Complete, 3 seeds & One newer target competes with old frequent prefixes \\
    Singleton & 8-way concurrent & Complete, 3 seeds & No shared prefix; overhead-only negative control \\
    Queue sharing & 16-way concurrent & Complete, 3 seeds & Preserve PEEK's current-queue strength and fairness proxy \\
    Strong hotspot switch & Sequential & Complete, 3 seeds & Nineteen old frequency-4 prefixes versus nineteen new frequency-2 prefixes \\
    80-request absence gap & Sequential & Complete, 3 seeds & Test request-window expiry and long-gap reuse \\
    Capacity sweep & Sequential & Complete, 3 seeds & 2,048/4,096/8,192 KV blocks \\
    Open-loop load/fairness & 4--16 req/s & Complete, 3 seeds & Saturation, tail latency, and hot/cold service equity \\
    Window sweep & Mixed & \todoresult & 64/128/256/512 admitted requests \\
    7B validation & Mixed & \todoresult & Primary performance model \\
    \bottomrule
  \end{tabularx}
\end{table*}

\section{Results}

\subsection{Initial six-policy mechanism check}

Before the paired three-seed matrix, we ran a single deterministic trace with
an exactly 3,456-token hot prefix, 24 unique pressure prefixes, and one hot
return. Table~\ref{tab:six-policy} establishes that PEEK's existing recency and
decay modes do not explain the H-PEEK mechanism: all non-history policies miss
the returning prefix, whereas both H-PEEK variants retain it. Because this was
one run and startup order was not randomized, it is mechanism evidence rather
than a performance estimate.

\begin{table*}[t]
  \centering
  \caption{Initial one-run six-policy pressure check at 4,096 KV blocks.}
  \label{tab:six-policy}
  \begin{tabular}{lrrr}
    \toprule
    Policy & Return hit/query tokens & Return compute tokens & Return TTFT (s) \\
    \midrule
    Vanilla vLLM & 0 / 3,465 & 3,465 & 0.3140 \\
    PEEK cluster & 0 / 3,465 & 3,465 & 0.3210 \\
    PEEK recency & 0 / 3,465 & 3,465 & 0.3140 \\
    PEEK decay & 0 / 3,465 & 3,465 & 0.3175 \\
    H-PEEK history-only & 3,456 / 3,465 & 9 & 0.0518 \\
    H-PEEK joint & 3,456 / 3,465 & 9 & 0.0526 \\
    \bottomrule
  \end{tabular}
\end{table*}

\subsection{Three-seed paired comparison}

Table~\ref{tab:main-results} reports arithmetic mean and sample standard
deviation across three deterministic prefix identities. Request throughput and
TTFT include the full workload. The positive trace contains one return request
per seed, reported separately in Table~\ref{tab:return}.

\begin{table*}[t]
  \centering
  \caption{End-to-end results at 4,096 KV blocks. Values are mean $\pm$ sample
  standard deviation over three seeds. Relative change is H-PEEK versus PEEK.}
  \label{tab:main-results}
  \begin{tabular}{llrrr}
    \toprule
    Workload & Metric & PEEK cluster & H-PEEK joint & Relative change \\
    \midrule
    Periodic reuse & Requests/s & $4.015\pm0.021$ & $4.118\pm0.014$ & $+2.57\%$ \\
                   & TTFT P95 (s) & $0.278\pm0.003$ & $0.278\pm0.001$ & $-0.01\%$ \\
                   & Prompt tokens/s & $12063.6\pm61.8$ & $12373.9\pm43.5$ & $+2.57\%$ \\
    \addlinespace
    Single-target drift & Requests/s & $8.999\pm0.055$ & $8.949\pm0.064$ & $-0.56\%$ \\
                        & TTFT P95 (s) & $0.274\pm0.002$ & $0.274\pm0.001$ & $-0.15\%$ \\
    \addlinespace
    Singleton & Requests/s & $12.908\pm0.051$ & $12.895\pm0.028$ & $-0.10\%$ \\
              & TTFT P95 (s) & $0.763\pm0.003$ & $0.764\pm0.001$ & $+0.06\%$ \\
              & Local cache-hit tokens & 0 & 0 & no change \\
    \addlinespace
    Queue sharing & Requests/s & $65.704\pm1.523$ & $65.179\pm0.964$ & $-0.80\%$ \\
                  & TTFT P95 (s) & $0.265\pm0.002$ & $0.263\pm0.005$ & $-1.01\%$ \\
                  & Local cache-hit tokens & 63,488 & 63,488 & no change \\
    \bottomrule
  \end{tabular}
\end{table*}

\begin{table}[t]
  \centering
  \caption{Periodic-reuse mechanism result.}
  \label{tab:return}
  \begin{tabular}{lrr}
    \toprule
    Metric & PEEK & H-PEEK \\
    \midrule
    Return TTFT (ms) & $272.5\pm0.8$ & $50.9\pm2.0$ \\
    Return TTFT change & --- & $-81.3\%$ \\
    Workload hit tokens & 8,976 & 11,968 \\
    Workload compute tokens & 78,155 & 75,163 \\
    Extra retained prefix & 0 & 2,992 tokens \\
    \bottomrule
  \end{tabular}
\end{table}

\subsection{Strong hotspot-switch experiment}

Both sides of the strong-switch experiment are complete. H-PEEK does not reduce
new-set cache hits under the default scan cap, but it adds measurable CPU-path
overhead. The free-queue scan cap limits global LFU-like reordering and thereby
preserves substantial recency behavior even when old prefixes have larger
historical counts.

\begin{table}[t]
  \centering
  \caption{Strong hotspot-switch results. Values are mean $\pm$ sample
  standard deviation over three seeds.}
  \label{tab:strong-switch-progress}
  \begin{tabular}{lrr}
    \toprule
    Metric & PEEK & H-PEEK \\
    \midrule
    Full workload requests/s & $9.171\pm0.091$ & $9.080\pm0.066$ \\
    Full workload throughput change & --- & $-1.00\%$ \\
    Full workload TTFT P95 (s) & $0.270\pm0.001$ & $0.272\pm0.001$ \\
    New-set return requests/s & $21.887\pm0.072$ & $21.140\pm0.334$ \\
    New-set return throughput change & --- & $-3.41\%$ \\
    New-set return TTFT P95 (s) & $0.049\pm0.002$ & $0.051\pm0.002$ \\
    New-set return TTFT P95 change & --- & $+4.10\%$ \\
    New-set return hit tokens & 56,848 & 56,848 \\
    New-set return compute tokens & $262\pm5$ & $262\pm5$ \\
    \bottomrule
  \end{tabular}
\end{table}

\subsection{History-evidence gate ablation}

The first H-PEEK ablation ignores historical frequency one: a block must occur
in at least two admitted requests before its history term contributes to the
protection score. The default configuration is denoted H-PEEK-$m{=}1$ and the
gated configuration H-PEEK-$m{=}2$. Both use a 256-request window and the same
PEEK cLPM scheduler. The gate preserves the positive periodic-reuse outcome:
the aggregate cached-token count is 35,904 for both history variants, compared
with 26,928 for PEEK on the same trace. It also reduces the stale-history
workload's CPU-side history ranking substantially.

\begin{table*}[t]
  \centering
  \small
  \caption{Minimum historical-frequency ablation at 4,096 KV blocks. Values are
  mean $\pm$ sample standard deviation over three seeds. TTFT values are in ms.}
  \label{tab:min-frequency}
  \begin{tabular}{llrrr}
    \toprule
    Workload & Metric & PEEK & H-PEEK-$m{=}1$ & H-PEEK-$m{=}2$ \\
    \midrule
    Periodic reuse & Requests/s & $4.015\pm0.021$ & $4.171\pm0.013$ & $4.182\pm0.010$ \\
                   & Return TTFT & $272.5\pm0.8$ & $46.5\pm3.2$ & $47.1\pm0.7$ \\
    Stale-history & Requests/s & $8.999\pm0.055$ & $9.147\pm0.131$ & $9.082\pm0.026$ \\
                  & Return TTFT & $46.9\pm4.7$ & $46.4\pm4.3$ & $49.4\pm1.3$ \\
    Singleton & Requests/s & $12.908\pm0.051$ & $12.876\pm0.045$ & $12.907\pm0.053$ \\
              & TTFT P95 & $763.4\pm2.8$ & $666.3\pm84.5$ & $713.2\pm84.8$ \\
    \bottomrule
  \end{tabular}
\end{table*}

\begin{table}[t]
  \centering
  \caption{Matched H-PEEK profile for the three-workload ablation.}
  \label{tab:min-frequency-overhead}
  \begin{tabular}{lrrr}
    \toprule
    Counter & $m{=}1$ & $m{=}2$ & Change \\
    \midrule
    Blocks scanned & 90,462 & 56,484 & $-37.6\%$ \\
    Eviction time (ms) & 346.5 & 192.5 & $-44.4\%$ \\
    History lookups & 53,451 & 19,486 & $-63.6\%$ \\
    Lookup time (ms) & 99.5 & 38.5 & $-61.3\%$ \\
    Changed priorities & 53,451 & 11,307 & $-78.8\%$ \\
    Observe time (ms) & 177.5 & 179.6 & $+1.1\%$ \\
    \bottomrule
  \end{tabular}
\end{table}

The gate is therefore a promising overhead optimization, but the current
three-seed latency differences are noisy and do not establish a universal
performance gain. The next experiment should test whether the same reduction
holds under open-loop load and whether a threshold of two misses legitimate
low-frequency reuse.

\subsection{History-window sensitivity}

The window-sensitivity run uses an 80-request absence gap. PEEK misses the
returning prefix in every seed and recomputes all 3,007 prompt tokens. H-PEEK
with $W=64$ also misses in every seed, confirming exact expiration in admitted-
request units. With $W=256$, H-PEEK instead retains 2,992 tokens and computes
only 15, reducing mean return TTFT from $268.6\pm2.0$ to $43.5\pm6.6$ ms
(83.8\%). $W=512$ produces the same exact hit/compute outcome and a
$45.7\pm0.8$ ms return TTFT. Thus, extending the window beyond the reuse
distance adds no cache-hit benefit on this trace.

\begin{table}[t]
  \centering
  \caption{History-window sensitivity after an 80-request absence gap. Values
  are means over three seeds unless shown with standard deviation.}
  \label{tab:window-progress}
  \begin{tabular}{lrrrrr}
    \toprule
    Metric & PEEK & $W=32$ & $W=64$ & $W=256$ & $W=512$ \\
    \midrule
    Return hit tokens & 0 & 0 & 0 & 2,992 & 2,992 \\
    Return compute tokens & 3,007 & 3,007 & 3,007 & 15 & 15 \\
    Return TTFT (ms) & $268.6\pm2.0$ & $267.5\pm2.8$ & $268.9\pm1.3$ & $43.5\pm6.6$ & $45.7\pm0.8$ \\
    Workload requests/s & $9.072\pm0.056$ & $9.145\pm0.101$ & $9.074\pm0.026$ & $9.290\pm0.069$ & $9.273\pm0.003$ \\
    \bottomrule
  \end{tabular}
\end{table}

\subsection{KV-capacity sweep in progress}

The paired 2,048-block runs are complete. Under this stronger pressure, the
periodic hot prefix is absent on PEEK return, as at 4,096 blocks. H-PEEK retains
2,992 tokens, reduces return TTFT by 81.8\%, and improves full-workload
throughput by 2.88\%. At 8,192 blocks, both policies already retain all 2,992
reusable tokens. H-PEEK then provides no hit benefit and reduces throughput by
0.63\%, isolating its low-pressure management overhead.

\begin{table}[t]
  \centering
  \caption{Periodic-reuse capacity sweep. Return hit tokens and TTFT are means
  over three seeds; dashes denote unfinished paired runs.}
  \label{tab:capacity-progress}
  \begin{tabular}{rrrll}
    \toprule
    Blocks & Policy & Return hits & Return TTFT (ms) & Requests/s \\
    \midrule
    2,048 & PEEK & 0 & $266.1\pm2.1$ & $4.057\pm0.015$ \\
          & H-PEEK & 2,992 & $48.5\pm1.1$ & $4.173\pm0.007$ \\
    4,096 & PEEK & 0 & $272.5\pm0.8$ & $4.015\pm0.021$ \\
          & H-PEEK & 2,992 & $50.9\pm2.0$ & $4.118\pm0.014$ \\
    8,192 & PEEK & 2,992 & $47.0\pm0.7$ & $4.202\pm0.011$ \\
          & H-PEEK & 2,992 & $48.2\pm0.8$ & $4.176\pm0.011$ \\
    \bottomrule
  \end{tabular}
\end{table}

\subsection{Open-loop offered load and fairness}

Each run contains
64 requests in a randomized 50/50 mixture: hot requests share a 2,048-token
prefix, while cold requests have unique 2,048-token prefixes. Inter-arrival
times are deterministic at the configured rate, and results are means over
three prefix-identity/order seeds. Both policies sustain the offered rate
through 8 requests/s, approach saturation at 12 requests/s, and saturate near
12 requests/s when offered 16. The client-side dispatch-lag P95 is at most
15.5 ms, so the observed multi-second tail at overload is not explained by the
load generator. Prefix-cache state is reset between runs. The history tracker
remains live within each policy server process; disjoint per-run Prefix IDs
prevent cross-run reuse, while retaining steady-state bounded-window overhead.

\begin{table*}[t]
  \centering
  \small
  \caption{Open-loop latency at 4,096 KV blocks. Latencies are seconds;
  values are mean $\pm$ sample standard deviation over three seeds.}
  \label{tab:open-loop-latency}
  \begin{tabular}{rrlrrrrr}
    \toprule
    Offered & Policy & Throughput & TTFT P50 & TTFT P95 & TTFT P99 & $\Delta$ throughput & $\Delta$ P95 \\
    \midrule
    4 & PEEK   & $4.021\pm0.024$ & $0.192\pm0.000$ & $0.198\pm0.001$ & $0.218\pm0.030$ & --- & --- \\
      & H-PEEK & $4.021\pm0.023$ & $0.192\pm0.000$ & $0.197\pm0.001$ & $0.215\pm0.029$ & $+0.02\%$ & $-0.15\%$ \\
    8 & PEEK   & $7.968\pm0.090$ & $0.200\pm0.006$ & $0.375\pm0.076$ & $0.458\pm0.098$ & --- & --- \\
      & H-PEEK & $7.967\pm0.092$ & $0.199\pm0.003$ & $0.381\pm0.073$ & $0.459\pm0.097$ & $-0.01\%$ & $+1.80\%$ \\
    12 & PEEK   & $11.426\pm0.162$ & $0.388\pm0.020$ & $0.901\pm0.375$ & $1.038\pm0.461$ & --- & --- \\
       & H-PEEK & $11.406\pm0.144$ & $0.394\pm0.020$ & $0.960\pm0.467$ & $1.048\pm0.456$ & $-0.18\%$ & $+6.57\%$ \\
    16 & PEEK   & $12.020\pm0.026$ & $0.494\pm0.011$ & $2.089\pm0.387$ & $2.671\pm0.734$ & --- & --- \\
       & H-PEEK & $11.994\pm0.009$ & $0.496\pm0.012$ & $2.100\pm0.375$ & $2.679\pm0.725$ & $-0.21\%$ & $+0.51\%$ \\
    \bottomrule
  \end{tabular}
\end{table*}

\begin{table*}[t]
  \centering
  \small
  \caption{Open-loop class latency and fairness. TTFT values are seconds.}
  \label{tab:open-loop-fairness}
  \begin{tabular}{rrlrrrr}
    \toprule
    Offered & Policy & Hot mean TTFT & Cold mean TTFT & Jain equity & TTFT $>1$s & Hit tokens \\
    \midrule
    4 & PEEK   & $0.046\pm0.001$ & $0.196\pm0.001$ & $0.722\pm0.003$ & 0 & 190,464 \\
      & H-PEEK & $0.046\pm0.001$ & $0.196\pm0.001$ & $0.721\pm0.004$ & 0 & 190,464 \\
    8 & PEEK   & $0.103\pm0.004$ & $0.263\pm0.011$ & $0.840\pm0.011$ & 0 & 190,464 \\
      & H-PEEK & $0.105\pm0.003$ & $0.266\pm0.014$ & $0.842\pm0.014$ & 0 & 190,464 \\
    12 & PEEK   & $0.302\pm0.010$ & $0.538\pm0.099$ & $0.927\pm0.036$ & $2.1\%\pm3.6\%$ & 190,464 \\
       & H-PEEK & $0.306\pm0.012$ & $0.544\pm0.096$ & $0.928\pm0.033$ & $2.1\%\pm3.6\%$ & 190,464 \\
    16 & PEEK   & $0.423\pm0.021$ & $1.145\pm0.109$ & $0.826\pm0.012$ & $23.4\%\pm1.6\%$ & 190,464 \\
       & H-PEEK & $0.428\pm0.018$ & $1.150\pm0.103$ & $0.827\pm0.012$ & $23.4\%\pm1.6\%$ & 190,464 \\
    \bottomrule
  \end{tabular}
\end{table*}

The two policies have identical cache-hit and recompute-token counts at every
load point. H-PEEK therefore provides no caching opportunity on this workload;
the paired data isolate overhead. Its throughput change ranges from $+0.02\%$
to $-0.21\%$. The largest mean P95 increase is 6.57\% at 12 requests/s, but
the across-seed standard deviation is 0.467 s and P99 rises only 0.96\%; more
repetitions are required before treating this noisy P95 difference as a stable
regression. Jain equity is computed over inverse mean hot/cold TTFT, so 1.0
denotes equal class-level service rate; neither fairness nor the one-second SLO
fraction changes materially.

The minimum-evidence gate also preserves the open-loop cache outcome: all four
offered loads have the same 190,464 cached prompt tokens for PEEK, H-PEEK-$m{=}1$,
and H-PEEK-$m{=}2$. Relative to the default history policy, the gate changes
throughput by at most $+0.04\%$ and keeps the 16-request/s P95 within 0.3\%.

\begin{table*}[t]
  \centering
  \small
  \caption{Open-loop minimum-evidence gate. Values are mean $\pm$ sample
  standard deviation over three seeds; TTFT P95 is in seconds.}
  \label{tab:open-loop-gate}
  \begin{tabular}{rllrr}
    \toprule
    Offered & Policy & Throughput & TTFT P95 & Jain equity \\
    \midrule
    4  & H-PEEK-$m{=}1$ & $4.021\pm0.023$ & $0.197\pm0.001$ & $0.721\pm0.004$ \\
       & H-PEEK-$m{=}2$ & $4.022\pm0.022$ & $0.198\pm0.002$ & $0.720\pm0.002$ \\
    8  & H-PEEK-$m{=}1$ & $7.967\pm0.092$ & $0.381\pm0.073$ & $0.842\pm0.014$ \\
       & H-PEEK-$m{=}2$ & $7.968\pm0.093$ & $0.383\pm0.075$ & $0.842\pm0.017$ \\
    12 & H-PEEK-$m{=}1$ & $11.406\pm0.144$ & $0.960\pm0.467$ & $0.928\pm0.033$ \\
       & H-PEEK-$m{=}2$ & $11.409\pm0.126$ & $0.906\pm0.366$ & $0.927\pm0.034$ \\
    16 & H-PEEK-$m{=}1$ & $11.994\pm0.009$ & $2.100\pm0.375$ & $0.827\pm0.012$ \\
       & H-PEEK-$m{=}2$ & $12.000\pm0.024$ & $2.094\pm0.376$ & $0.824\pm0.009$ \\
    \bottomrule
  \end{tabular}
\end{table*}

In the matched open-loop profile, the gate reduces blocks scanned by 28.3%,
eviction time by 39.9%, history lookups by 89.9%, and lookup time by 87.5%.
This is consistent with the three-workload profile and indicates that the gate
reduces real scheduler work rather than only changing an offline score.

\subsection{Internal overhead}

The profile in Table~\ref{tab:overhead} covers 720 admitted requests across the
four completed workloads. End-to-end singleton throughput remains nearly
unchanged because GPU execution dominates this small experiment, but H-PEEK's
CPU-side scan cost is materially larger and should not be hidden.

\begin{table}[t]
  \centering
  \caption{Aggregate scheduler/eviction profile.}
  \label{tab:overhead}
  \begin{tabular}{lrr}
    \toprule
    Counter & PEEK & H-PEEK \\
    \midrule
    Eviction calls & 1,016 & 1,014 \\
    Blocks scanned & 38,478 & 91,446 \\
    Eviction time (ms) & 101.6 & 341.1 \\
    History observations & 0 & 720 \\
    History update time (ms) & 0 & 205.3 \\
    History lookups & 0 & 53,953 \\
    History lookup time (ms) & 0 & 90.1 \\
    Lookup time/candidate ($\mu$s) & --- & 1.67 \\
    Window key references & 0 & 33,308 \\
    \bottomrule
  \end{tabular}
\end{table}

The open-loop profile provides a second overhead measurement over 768 admitted
requests. Both policies execute exactly 1,164 eviction calls. H-PEEK scans
40.1\% more blocks and spends 81.6\% more profiled time in eviction; its
history admission update costs 0.287 ms/request and each historical lookup
costs 1.64 $\mu$s on average.

\begin{table}[t]
  \centering
  \caption{Open-loop aggregate CPU-path profile.}
  \label{tab:open-loop-overhead}
  \begin{tabular}{lrr}
    \toprule
    Counter & PEEK & H-PEEK \\
    \midrule
    Admitted requests & 768 & 768 \\
    Eviction calls & 1,164 & 1,164 \\
    Blocks scanned & 51,152 & 71,686 \\
    Eviction time (ms) & 113.8 & 206.7 \\
    History update time (ms) & 0 & 220.1 \\
    History lookups & 0 & 22,445 \\
    History lookup time (ms) & 0 & 36.9 \\
    \bottomrule
  \end{tabular}
\end{table}

\section{Interpretation and Threats to Validity}

The periodic result demonstrates the intended mechanism: PEEK cannot observe a
future request while its prefix is absent from the queue, whereas admitted-
request frequency retains the prefix. The singleton and queue-sharing results
show no large end-to-end regression at this scale. However, the 2.38$\times$
scan-volume increase and 3.36$\times$ aggregate eviction-time increase are real
costs. The strong-switch experiment does not show a cache-hit regression under
the default scan cap, but its new-set return phase is 3.41\% slower and has
4.10\% higher TTFT P95. Thus, the immediate observed failure mode is CPU scan
overhead rather than stale-history cache pollution.

The study currently uses one model size, one GPU, one cache block size, three
prefix-identity seeds, and synthetic traces. Server startup is excluded, but
request execution is not repeated enough for publication-grade tail-latency
claims. A complete evaluation must include policy-order randomization, more
seeds, confidence intervals, capacity and history-window sweeps, 7B validation,
and the full PEEK paper workloads. The open-loop process resets KV-cache state
but not tracker history between its 12 disjoint-prefix runs. This preserves a
steady-state 256-request tracker footprint, but a publication run should add an
explicit tracker-reset hook or restart the server per run to separate cold and
steady-state overhead.

\section{Planned Ablations}

If strong-switch pollution is observed, we will evaluate bounded historical
scores, minimum-frequency activation, and request-index aging. Request-index
aging is based on admitted-request distance rather than wall-clock block access
time, keeping it conceptually distinct from PEEK's existing recency and decay
modes. For overhead, we will test a fast path that bypasses history-aware scans
when history supplies no rank discrimination, such as an all-singleton trace.

\section{Reproducibility}

Raw JSON results, engine profiles, environment pins, and the patch relative to
the official PEEK commit are stored in the artifact repository. The benchmark
driver records exact tokenized prefix lengths and obtains metric deltas from
vLLM's Prometheus endpoint. Administrative cache reset endpoints are enabled
only on the loopback interface for experimental isolation.

\end{document}
